\section{Exact Selection Across Structural Families}
\label{app:family-selection}

The primary decoder fixes a feasible family and reports its constraint tension.
Multiple inheritance raises a different question: does the
available evidence justify changing the family itself?  We formalize that
question for nested families.  Let $\mathcal F_0$ contain single-parent
forests and let $\mathcal F_B\supseteq\mathcal F_0$ contain acyclic
taxonomies with at most two parents per child and at most $B$ second-parent
slots.  Write $A_s(y)$ for the unregularized additive evidence and
$M(y)$ for the number of used second-parent slots.  Define
\begin{align}
 P_B(s)&=\max_{y\in\mathcal F_B} A_s(y),\\
 R_B(s)&=P_B(s)-P_0(s),\label{eq:relaxation-value}\\
 L_B(s)&=\max_{y\in\mathcal L_B} A_s(y),\\
 T_B(s)&=L_B(s)-P_B(s).\label{eq:family-tension}
\end{align}
where $\mathcal L_B$ drops global acyclicity and budget coupling.
$R_B$ is the \emph{relaxation value}: neural score recovered by enlarging
the forest family. $T_B$ is the \emph{constraint tension}: score difference required
to enforce the declared global constraints.  They answer different questions.

Since $\mathcal F_0\subseteq\mathcal F_B\subseteq\mathcal L_B$,
$T_B,R_B\geq0$ and the identity
$L_B(s)-P_0(s)=T_B(s)+R_B(s)$ separates local evidence, global constraint
tension, and the score value of multiple inheritance.

Rather than report only two dev-selected thresholds, consider the complete
regularization path
\begin{align}
 y_\lambda&\in\argmax_{y\in\mathcal F_B}
 \{A_s(y)-\lambda M(y)\},\\
 V(\lambda)&=A_s(y_\lambda)-\lambda M(y_\lambda).
 \label{eq:family-path}
\end{align}
Positive $\lambda$ penalizes additional parents; negative $\lambda$ is an
explicit structural subsidy and is reported as such.

\begin{proposition}[Monotone structural path]
For $\lambda_1<\lambda_2$ and arbitrary optimizers at the two values,
$M(y_{\lambda_1})\geq M(y_{\lambda_2})$.  Moreover, $V(\lambda)$ is
convex and piecewise linear with an active slope $-M(y_\lambda)$.
\end{proposition}
\begin{proof}
Optimality at the two penalties gives
$A(y_1)-\lambda_1M(y_1)\geq A(y_2)-\lambda_1M(y_2)$ and
$A(y_2)-\lambda_2M(y_2)\geq A(y_1)-\lambda_2M(y_1)$.
Adding yields
$(\lambda_2-\lambda_1)(M(y_1)-M(y_2))\geq0$.
The value is the pointwise maximum of finitely many affine functions.
\end{proof}

The implementation obtains exact breakpoints from the cardinality frontier
$P_m=\max_{y:M(y)=m}A_s(y)$.  Cardinality $m$ is optimal precisely on
\begin{equation}
 \max_{j>m}\frac{P_j-P_m}{j-m}
 \;\leq\lambda\leq\;
 \min_{j<m}\frac{P_m-P_j}{m-j},
 \label{eq:penalty-interval}
\end{equation}
when this interval is nonempty.  These intervals certify that a reported
forest/DAG decision is not an artifact of a single threshold grid point.

\paragraph{Family-stability radius.}
Let $q(y)$ be the selected edge/ROOT indicator vector and let $y^*$ be the
current optimum.  Let $\gamma_{\rm fam}$ be its objective gap to the best
solution in the other family (forest versus any solution using a second
parent), and let
\begin{equation}
 D=\max_{y\text{ in other family}}\|q(y^*)-q(y)\|_1.
\end{equation}
Then
\begin{equation}
 \rho_{\rm fam}=\gamma_{\rm fam}/D
 \label{eq:family-radius}
\end{equation}
is a conservative family-stability radius: every perturbation
$\|\Delta s\|_\infty<\rho_{\rm fam}$ preserves the selected family.
Indeed, the score difference against every other-family solution decreases
by at most $\|\Delta s\|_\infty D<\gamma_{\rm fam}$.
Forced exclusion of a selected edge and forced inclusion of an unselected
second-parent decision similarly give edge-necessity and inclusion gaps.

\paragraph{A tractable ordered case.}
If every candidate edge respects a fixed strict topological order, any subset
is acyclic.  For each child, compute its best zero/one-parent value $b_c$ and
best two-parent value $d_c$, then select the at most $B$ largest positive
margins $d_c-b_c$.  This is globally exact because the budget is the only
remaining cross-child constraint.  It costs
$O(|E|+|V|\log|V|)$ for additive edge utilities, or
$O(\sum_c k_c^2+|V|\log|V|)$ with arbitrary parent-pair potentials.
The general cyclic candidate case retains the ILP decoder.

\paragraph{Evidence model and failure attribution.}
Raw single-parent ranking log-potentials are not variable-cardinality
probabilities.  Our revised pipeline therefore fits edge log-odds on WordNet
train, a child-level cardinality head, and a conditional parent-pair head.
The pair head ranks pairs after the first-parent evidence is known, while the
cardinality head decides whether the second slot is warranted.  WordNet dev
is reserved for ROOT calibration and $\lambda$ selection; the hard cap is a
predeclared structural limit rather than gold prevalence.

Retrieval remains an independent ceiling.  Among the
\MpOracleGoldChildren{} Macro4 multi-parent children, both gold parents occur
in the frozen top-20 pool for only
\MpOracleCompleteTwenty{} (\MpOracleCompleteTwentyRate\%) and in top-50 for
\MpOracleCompleteFifty{} (\MpOracleCompleteFiftyRate\%).  Food contributes
only \MpOracleFoodTwenty{} complete top-20 cases out of
\MpOracleFoodGold{}.  We consequently report four stages separately:
candidate-set coverage, cardinality detection, conditional second-parent
ranking, and exact globally feasible decoding.

\paragraph{Sealed semantic-pair follow-up.}
After inspecting the hybrid-retrieval and pair-oracle outcomes, we froze a
separate ten-seed follow-up; it cannot alter the primary branching or
allocation claims.  It reuses the coverage-preserving hybrid top-50 pools and
frozen 768-dimensional edge representations, requiring
\PairTransformerForwards{} additional transformer forwards.  For each child,
the scorer enumerates unordered pairs within the top-$q=20$ edge frontier and
uses the symmetric representation
\begin{equation}
 [h_i+h_j,\ |h_i-h_j|,\ h_i\odot h_j,\ \phi_{ij}],
\end{equation}
where $h_i$ is a learned transformation of frozen edge evidence and $\phi_{ij}$
contains symmetric scalar rank and score features.  The exact decoder sees
only the top-ranked pair ($M=1$), so the neural component must select the pair
and the symbolic component only enforces global DAG feasibility.

\begin{table}[h]
\centering
\scriptsize
\setlength{\tabcolsep}{3pt}
\begin{tabular}{lrrrrr}
\toprule
& \multicolumn{2}{c}{\ef} & \multicolumn{2}{c}{\af} & Pair R@1\\
\cmidrule(lr){2-3}\cmidrule(lr){4-5}
Scope & Hybrid & Sem.+margin & Hybrid & Sem.+margin & Scalar$\to$sem.\\
\midrule
Macro4 & \PairMacroFourBaselineEdge & \textbf{\PairMacroFourSemanticEdge} &
\PairMacroFourBaselineAncestor & \textbf{\PairMacroFourSemanticAncestor} &
\PairMacroFourScalarPairRecallOne$\rightarrow$\textbf{\PairMacroFourSemanticPairRecallOne}\\
Macro6 & \PairMacroSixBaselineEdge & \textbf{\PairMacroSixSemanticEdge} &
\PairMacroSixBaselineAncestor & \textbf{\PairMacroSixSemanticAncestor} &
\PairMacroSixScalarPairRecallOne$\rightarrow$\textbf{\PairMacroSixSemanticPairRecallOne}\\
\bottomrule
\end{tabular}
\caption{Post-outcome semantic parent-pair follow-up (mean percent, seeds
67--76). Pair R@1 is conditional on both gold parents occurring in the frozen
top-50 pool; its matched scalar control uses the same candidates and decoder.}
\label{tab:semantic-pair-followup}
\end{table}

On Macro4, semantic shortlisting improves \ef{} by
\PairMacroFourEdgeDelta{} points (95\% paired-seed bootstrap CI
[\PairMacroFourEdgeLow, \PairMacroFourEdgeHigh], exact sign-flip
$p=\PairMacroFourEdgeP$) and \af{} by \PairMacroFourAncestorDelta{} points
([\PairMacroFourAncestorLow, \PairMacroFourAncestorHigh],
$p=\PairMacroFourAncestorP$).  The edge effect is positive on
\PairMacroFourEdgeWins/\PairSeeds{} seeds.  Against the matched scalar pair
score, semantic pair recall@1 gains \PairMacroFourPairRecallDelta{} points
([\PairMacroFourPairRecallLow, \PairMacroFourPairRecallHigh],
$p=\PairMacroFourPairRecallP$).  Supporting Macro6 analysis shows
\PairMacroSixEdgeDelta/\PairMacroSixAncestorDelta{} \ef/\af{} gains
($p=\PairMacroSixEdgeP/\PairMacroSixAncestorP$) and a
\PairMacroSixPairRecallDelta-point pair-recall gain
($p=\PairMacroSixPairRecallP$).  Forest outputs are exactly unchanged, which
localizes these gains to semantic restriction of the multi-parent feasible
set.  Exact parent-set recovery is observed too rarely to support a positive
claim; the follow-up therefore establishes improved pair ranking and global
hierarchy reconstruction, not solved multiple inheritance.

\paragraph{Confirmatory joint-set budget ablation.}
We finally evaluate all predeclared pair frontiers with the joint
ROOT/singleton/pair scorer on a separate sealed set of
\JointSetSeeds{} seeds. A fixed frontier of $q=3$ scores only
\JointQThreePairs{} unordered pairs per child, compared with
\JointQFivePairs{} at $q=5$, \JointQTenPairs{} at $q=10$, and
\JointAdaptivePairs{} under uncertainty-adaptive allocation.

\begin{table}[h]
\centering
\scriptsize
\setlength{\tabcolsep}{3pt}
\begin{tabular}{llrrrr}
\toprule
Retrieval & Frontier & Pairs/child & Exact-set F1 & \ef & \af\\
\midrule
Hybrid40 & $q=3$ & \JointQThreePairs & \textbf{\JointHybridQThreeExact} &
\textbf{\JointHybridQThreeEdge} & \textbf{\JointHybridQThreeAncestor}\\
& $q=5$ & \JointQFivePairs & \JointHybridQFiveExact &
\JointHybridQFiveEdge & \JointHybridQFiveAncestor\\
& $q=10$ & \JointQTenPairs & \JointHybridQTenExact &
\JointHybridQTenEdge & \JointHybridQTenAncestor\\
& adaptive & \JointAdaptivePairs & \JointHybridAdaptiveExact &
\JointHybridAdaptiveEdge & \JointHybridAdaptiveAncestor\\
\midrule
Learned50 & $q=3$ & \JointQThreePairs & \textbf{\JointLearnedQThreeExact} &
\textbf{\JointLearnedQThreeEdge} & \textbf{\JointLearnedQThreeAncestor}\\
& $q=5$ & \JointQFivePairs & \JointLearnedQFiveExact &
\JointLearnedQFiveEdge & \JointLearnedQFiveAncestor\\
& $q=10$ & \JointQTenPairs & \JointLearnedQTenExact &
\JointLearnedQTenEdge & \JointLearnedQTenAncestor\\
& adaptive & \JointAdaptivePairs & \JointLearnedAdaptiveExact &
\JointLearnedAdaptiveEdge & \JointLearnedAdaptiveAncestor\\
\bottomrule
\end{tabular}
\caption{Confirmatory Macro4 joint-set budget ablation (mean percent, seeds
77--86). Bold marks the best mean within each retrieval pool; all rows retain
exact DAG feasibility.}
\label{tab:joint-set-budget}
\end{table}

The smallest frontier has the best mean exact-set, edge, and ancestor F1 in
both retrieval pools while reducing pair evaluation by
\JointQThreeSavingAdaptive\% relative to adaptive allocation,
\JointQThreeSavingQFive\% relative to $q=5$, and
\JointQThreeSavingQTen\% relative to $q=10$. We interpret this as a
descriptive efficiency result: wider pair enumeration provides no observed
benefit once the joint scorer has identified the strongest local frontier.
The preregistered exact-set contrast between Learned50+adaptive and
Hybrid40+$q=5$ is \JointPrimaryCandidateExact{} versus
\JointPrimaryReferenceExact{} F1 ($\Delta=\JointPrimaryExactDelta$ points,
Holm $p=\JointPrimaryExactP$), so it does not support a claim of improved
exact parent-set recovery.
